\subsection{梯度场的刻画. de Rham 定理}
\label{sec:chap4-gradient-fields}

\renewcommand{\thestatement}{2.\arabic{statement}}
\setcounter{statement}{0}
\newcounter{chapterfoursectiontworemark}
\renewcommand{\thechapterfoursectiontworemark}{2.\arabic{chapterfoursectiontworemark}}

本节的目标是在所有向量场中给出梯度场的一个有用刻画.

\begin{Definition}
\label{def:chap4-annihilator}
设 $E$ 是 Banach 空间, 其对偶空间为 $E'$. 对 $E$ 的任意子集 $A$, 定义 $A$ 的正交补为 $E'$ 的线性子空间
\[
 A^\perp=\{\varphi\in E',\ \forall x\in A,\ \langle\varphi,x\rangle_{E',E}=0\}.
\]
\end{Definition}

当 $E$ 是 Hilbert 空间时, 利用 Riesz 表示定理容易看出, $A^\perp$ 也就是相对于 $E$ 的内积而言的 $A$ 的正交补.

\par\noindent\refstepcounter{chapterfoursectiontworemark}\textbf{Remark~\thechapterfoursectiontworemark:}\label{rem:chap4-annihilator-properties}\quad
\begin{itemize}
 \item $A$ 的正交补是 $E'$ 的闭线性子空间. 当 $E$ 自反时, $A^{\perp\perp}$ 可以与 $E$ 的一个线性子空间等同. 此时它是包含 $A$ 的最小闭线性子空间; 即 $A^{\perp\perp}=\overline{\operatorname{Span}A}$. 因此, 对 $E$ 的任意向量子空间 $A$, 有
 \[
  A\text{ 在 }E\text{ 中闭}\quad\Longleftrightarrow\quad A^{\perp\perp}=A.
 \]
 \item 若 $A\subset B$, 则 $B^\perp\subset A^\perp$.
\end{itemize}

\begin{Theorem}
\label{thm:chap4-gradient-image-closed}
设 $\Omega$ 是 $\mathbb R^d$ 中的 Lipschitz 区域.
\begin{itemize}
 \item 若 $\Omega$ 有界, 则梯度算子的像
 \[
  \nabla\bigl(L^2(\Omega)\bigr)
  =\{f\in(H^{-1}(\Omega))^d,\ \exists p\in L^2(\Omega),\ f=\nabla p\}
 \]
 在 $(H^{-1}(\Omega))^d$ 中闭.
 \item 若 $\Omega$ 无界, 则集合
 \[
  \nabla\bigl(L_{\mathrm{loc}}^2(\Omega)\bigr)\cap(H^{-1}(\Omega))^d
  =\{f\in(H^{-1}(\Omega))^d,\ \exists p\in L_{\mathrm{loc}}^2(\Omega),\ f=\nabla p\}
 \]
 在 $(H^{-1}(\Omega))^d$ 中闭.
\end{itemize}
\end{Theorem}

注意, 在无界情形下, $L_{\mathrm{loc}}^2(\Omega)$ 中函数的梯度 (按分布意义) 一般并不是 $(H^{-1}(\Omega))^d$ 的元素.

\begin{Proof}[Theorem~\ref{thm:chap4-gradient-image-closed}]
只需考虑连通区域的情形. 事实上, 若 $\Omega$ 不连通, 则在它的每个连通分支上应用该定理即可得到结果.

\begin{itemize}
 \item 首先假设 $\Omega$ 有界. 设 $(\nabla p_n)_n$ 是 $(H^{-1}(\Omega))^d$ 中收敛的序列, 其中 $p_n\in L^2(\Omega)$. 由于 $\Omega$ 有界, 可以在不改变 $\nabla p_n$ 的情况下假设 $m(p_n)=0$.

 现在应用 Nečas 不等式 (Theorem~\ref{thm:chap4-necas-inequality}) 和 Poincaré 不等式 (Proposition~\ref{prop:chap4-hminusone-poincare}), 得到
 \[
  \|p_n-p_{n+k}\|_{L^2}
  \leq C\|\nabla(p_n-p_{n+k})\|_{H^{-1}}.
 \]
 由于 $(\nabla p_n)_n$ 是 $(H^{-1}(\Omega))^d$ 中的 Cauchy 序列, 推出 $(p_n)_n$ 是 $L^2(\Omega)$ 中的 Cauchy 序列. 因而存在 $p\in L^2(\Omega)$, 使 $n\to\infty$ 时 $p_n\to p$. 所以
 \[
  \nabla p_n\xrightarrow[n\to\infty]{}\nabla p
  \quad\text{in }(H^{-1}(\Omega))^d,
 \]
 结论得证.

 \item 在无界情形下, 从有界情形的结果推出结论. 为此, 考虑一列递增的正则、有界且连通的开集 $(\Omega_k)_k$, 其并集为 $\Omega$.

 现在设 $(f_n)_n$ 是 $(H^{-1}(\Omega))^d$ 中的元素序列, 可写成 $f_n=\nabla p_n$, 其中 $p_n\in L_{\mathrm{loc}}^2(\Omega)$, 并且该序列收敛到 $(H^{-1}(\Omega))^d$ 中的元素 $f$.

 对每个 $k\geq1$, $f$ 在 $\Omega_k$ 上的限制是良定义的 (把 $H_0^1(\Omega_k)$ 中的函数在整个 $\Omega$ 上作零延拓). 此外, 容易看出, 序列 $(f_n|_{\Omega_k})_n$ 在 $(H^{-1}(\Omega_k))^d$ 中收敛到 $f|_{\Omega_k}$. 因而, 由有界情形的结果, 已知 $f|_{\Omega_k}$ 是某个函数 $q_k\in L^2(\Omega_k)$ 的梯度.

 由于 $\Omega_k$ 连通, 函数 $q_k$ 在相差一个常数的意义下唯一. 选取 $q_k$, 使它在 $\Omega_1$ 上均值为零 ($\Omega_1$ 包含在所有 $\Omega_k$ 中).

 因此, 若 $k_1<k_2$, 则函数 $q_{k_1}-q_{k_2}$ 在 $\Omega_{k_1}$ 上为零. 事实上, 该函数的分布梯度为零 (因为 $f|_{\Omega_{k_1}}=\nabla q_{k_1}$ 且 $f|_{\Omega_{k_2}}=\nabla q_{k_2}$), 所以根据 Lemma~\ref{lem:chap2-zero-gradient-distribution-constant}, 它是常数. 然而, $q_{k_1}$ 和 $q_{k_2}$ 在 $\Omega_1$ 上的均值都为零, 因而这个常数必为零.

 因此, 函数 $(q_k)_k$ 在 $(\Omega_k)_k$ 上相容地构造, 这给出定义在整个 $\Omega$ 上的函数 $q$, 且确有 $q\in L_{\mathrm{loc}}^2(\Omega)$. 此外, 显然 $q$ 的梯度在 $(H^{-1}(\Omega))^d$ 中也等于 $f$.
\end{itemize}
\end{Proof}

\begin{Theorem}
\label{thm:chap4-de-rham-hminusone}
设 $\Omega$ 是 $\mathbb R^d$ 中边界紧的连通 Lipschitz 区域. 设 $f\in(H^{-1}(\Omega))^d$, 并且
\[
 \langle f,\varphi\rangle_{H^{-1},H_0^1}=0,
 \qquad\forall\varphi\in(H_0^1(\Omega))^d\text{ 满足 }\operatorname{div}\varphi=0.
\]
则存在 $p\in L_{\mathrm{loc}}^2(\Omega)$, 它在相差一个常数的意义下唯一, 使得 $f=\nabla p$.

此外, 若 $\Omega$ 有界, 则 $p\in L^2(\Omega)$.
\end{Theorem}

\begin{Proof}
令 $Y$ 是 $(H^{-1}(\Omega))^d$ 的如下子空间:
\[
 \begin{cases}
 Y=\{\nabla p,\ p\in L^2(\Omega)\},&\text{若 }\Omega\text{ 有界},\\
 Y=\{f\in(H^{-1}(\Omega))^d,\ \exists p\in L_{\mathrm{loc}}^2(\Omega),\ f=\nabla p\},&\text{若 }\Omega\text{ 无界}.
 \end{cases}
\]
再令 $Z$ 是 $(H_0^1(\Omega))^d$ 的子空间
\[
 Z=\{\varphi\in(H_0^1(\Omega))^d,\ \operatorname{div}\varphi=0\}.
\]

由 Definition~\ref{def:chap4-annihilator}, 要证明的定理断言: 若 $f\in Z^\perp$, 则 $f\in Y$. 然而, Theorem~\ref{thm:chap4-gradient-image-closed} 表明 $Y$ 是自反空间 $(H^{-1}(\Omega))^d$ 的闭子空间, 因而由 Remark~\ref{rem:chap4-annihilator-properties} 有 $Y^{\perp\perp}=Y$. 所以只需证明包含关系 $Y^\perp\subset Z$ (另一个包含关系显然成立).

设 $u\in Y^\perp\subset(H_0^1(\Omega))^d$. 按定义, 这意味着
\[
 \langle\nabla p,u\rangle_{H^{-1},H_0^1}=0,
 \qquad\forall p\in L^2(\Omega).
\]
另一方面, 对任意 $p\in L^2(\Omega)$, 按定义有
\[
 \langle\nabla p,u\rangle_{H^{-1},H_0^1}
 =-\int_\Omega p\operatorname{div}u\,\mathrm dx.
\]
因此已经证明, 任意 $u\in Y^\perp$ 都满足
\[
 \int_\Omega p(\operatorname{div}u)\,\mathrm dx=0,
 \qquad\forall p\in L^2(\Omega).
\]
在这个公式中取 $p=\operatorname{div}u$, 恰好得到 $\operatorname{div}u=0$, 因而 $u\in Z$.
\end{Proof}

现在可以证明 Theorem~\ref{thm:chap4-de-rham-hminusone} 的一个更强版本, 它在更弱的假设下给出相同结论. 事实上, 现在可以假设 $f$ 在 $(\mathcal D(\Omega))^d$ 中的无散度测试函数上为零, 而不再要求它在 $(H_0^1(\Omega))^d$ 的所有无散度函数上为零.

\begin{Theorem}{de Rham}
\label{thm:chap4-de-rham}
设 $\Omega$ 是 $\mathbb R^d$ 中连通、有界的 Lipschitz 区域. 设 $f\in(H^{-1}(\Omega))^d$, 并且对任意满足 $\operatorname{div}\varphi=0$ 的函数 $\varphi\in(\mathcal D(\Omega))^d$, 有
\[
 \langle f,\varphi\rangle_{H^{-1},H_0^1}=0.
\]
则存在唯一函数 $p\in L_0^2(\Omega)$, 使得 $f=\nabla p$.
\end{Theorem}

\begin{Proof}
\begin{itemize}
 \item 第一步:

 设 $(\Omega_n)_n$ 是一列包含于 $\Omega$ 的递增正则连通开集 (即 $\overline{\Omega_n}\subset\Omega_{n+1}$), 使 $\bigcup_n\Omega_n=\Omega$.

 对任意满足 $\operatorname{div}u_n=0$ 的 $u_n\in(H_0^1(\Omega_n))^d$ 以及任意 $\varepsilon>0$, 令
 \[
  u_{n,\varepsilon}=\overline{u_n}*\eta_\varepsilon.
 \]
 注意到, 当 $\varepsilon$ 足够小时, $u_{n,\varepsilon}$ 紧支撑于 $\Omega$, 因为 $u_n$ 支撑于 $\Omega_n$. 因此 $u_{n,\varepsilon}\in(\mathcal D(\Omega))^d$ 且 $\operatorname{div}u_{n,\varepsilon}=0$. 由假设,
 \[
  \langle f,u_n\rangle_{H^{-1},H_0^1}
  =\lim_{\varepsilon\to0}\langle f,u_{n,\varepsilon}\rangle_{H^{-1},H_0^1}=0.
 \]
 所以由 Theorem~\ref{thm:chap4-de-rham-hminusone}, $f|_{\Omega_n}$ 作为 $(H^{-1}(\Omega_n))^d$ 的元素, 是某个函数 $p_n\in L^2(\Omega_n)$ 的梯度. 此外, 对所有 $n$, $p_{n+1}$ 在 $\Omega_n$ 上的限制与 $p_n$ 具有相同梯度, 因而 $p_{n+1}-p_n$ 在 $\Omega_n$ 中是常函数; 可以选取 $p_{n+1}$, 使这个常数为零.

 总之, 函数序列 $(p_n)_n$ 定义了一个函数 $p\in L_{\mathrm{loc}}^2(\Omega)$, 满足
 \[
  f=\nabla p.
 \]

 \item 第二步:

 还需证明 $p$ 的确属于 $L^2(\Omega)$. 为此, 使用 Lemma~\ref{lem:chap3-lipschitz-locally-star-shaped}. 于是可以写成
 \[
  \Omega=\bigcup_{i=1}^k\omega_i,
 \]
 其中 $\omega_i$ 是有限多个星形开集. 只需证明对所有 $i$ 都有 $p\in L^2(\omega_i)$.

 现在置身于星形开集 $\omega_i$ 中; 为简化记号, 在平移之后假设它关于 $0$ 星形. 对任意满足 $\operatorname{div}u_i=0$ 的 $u_i\in(H_0^1(\omega_i))^d$ (把它在整个 $\mathbb R^d$ 上作零延拓), 以及任意 $\theta\in[0,1[$, 令 $u_{i,\theta}(x)=u_i(x/\theta)$. 它紧支撑于 $\omega_i$, 因而紧支撑于 $\Omega$, 且无散度. 再通过正则化, 引入
 \[
  u_{i,\theta,\varepsilon}=u_{i,\theta}*\eta_\varepsilon,
 \]
 它无散度, 且当 $\varepsilon>0$ 足够小时紧支撑于 $\Omega$, 所以假设给出
 \[
  \langle f,u_{i,\theta}\rangle_{H^{-1},H_0^1}
  =\lim_{\varepsilon\to0}\langle f,u_{i,\theta,\varepsilon}\rangle_{H^{-1},H_0^1}=0.
 \]
 接着容易证明, 当 $\theta\to1$ 时, $u_{i,\theta}$ 在 $(H_0^1(\Omega))^d$ 中收敛到 $u_i$; 最终得到
 \[
  \langle f,u_i\rangle_{H^{-1},H_0^1}
  =\lim_{\theta\to1}\langle f,u_{i,\theta}\rangle_{H^{-1},H_0^1}=0.
 \]
 对所有 $u_i\in(H_0^1(\omega_i))^d$ 都有这个结果, 因而可以应用 Theorem~\ref{thm:chap4-de-rham-hminusone}, 得知 $f=\nabla p_i$, 其中 $p_i\in L^2(\omega_i)$. 与第一步一样, 显然 $p-p_i$ 是常数, 并且可以选取这个常数为零. 因此已经证明
 \[
  p|_{\omega_i}\in L^2(\omega_i),
  \qquad\forall i\in\{1,\ldots,k\},
 \]
 从而 $p\in L^2(\Omega)$. 当然, 现在可以用 $p-m(p)$ 代替 $p$, 以保证 $p\in L_0^2(\Omega)$.
\end{itemize}
\end{Proof}

事实上, 还有一个由 de Rham 得到的更一般结果 (见\MBUnresolvedReference{bib:de-rham-general}{[49]}), 其推论之一如下 (另见\MBUnresolvedReference{bib:de-rham-related}{[111, 112, 124]}).

\begin{Theorem}
\label{thm:chap4-de-rham-distributions}
设 $\Omega$ 是 $\mathbb R^d$ 的任意开集, 且 $f\in(\mathcal D'(\Omega))^d$. 若对所有散度为零的 $\varphi\in(\mathcal D(\Omega))^d$, 都有
\[
 \langle f,\varphi\rangle_{\mathcal D',\mathcal D}=0,
\]
则存在 $\pi\in\mathcal D'(\Omega)$, 使得 $f=\nabla\pi$.
\end{Theorem}
